\documentclass[journal]{IEEEtran}

\usepackage{amsmath}
\usepackage{amssymb}
\interdisplaylinepenalty=2500

\hyphenation{op-tical net-works semi-conduc-tor}

\begin{document}

\title{Notification-Gated Reliability Unit Commitment:\\ Least-Distortionary Out-of-Market Interventions}

\author{Author~One,~\IEEEmembership{Member,~IEEE,}
        Author~Two,
        and~Author~Three,~\IEEEmembership{Senior~Member,~IEEE}%
\thanks{Affiliations and funding acknowledgments.}%
\thanks{Manuscript received XXXX; revised XXXX.}}

\markboth{IEEE Transactions on Power Systems}%
{Authors: Notification-Gated Reliability Unit Commitment}

\maketitle

\begin{abstract}
System operators routinely commit generation resources outside of the
energy market to ensure reliability, but existing reliability unit
commitment (RUC) formulations optimize over the full planning horizon
without distinguishing decisions that must be made immediately from
those that can be deferred to subsequent market processes.
This scope mismatch leads to systematic over-commitment, displacing
market-scheduled resources and distorting price signals.
We propose a notification-gated RUC formulation that uses physical
generator startup lead times to identify the minimal set of
non-deferrable commitment decisions at each operational stage.
The formulation is a mixed-integer second-order cone program (MISOCP)
that minimizes reliability intervention costs over the gating window
while enforcing worst-case dispatch feasibility under ellipsoidal
uncertainty via linear decision rules.
The gating framework generalizes across operational horizons ---
multi-day assessment, day-ahead RUC, and intra-day RUC are all
instances with different parameterizations --- and recovers standard
DARUC as a limiting case.
Numerical experiments on the RTS-GMLC test case demonstrate
significant reductions in out-of-market commitments while maintaining
system reliability under worst-case uncertainty.
\end{abstract}

\begin{IEEEkeywords}
Reliability unit commitment, robust optimization, market design,
notification time, linear decision rules.
\end{IEEEkeywords}

\IEEEpeerreviewmaketitle


%% ===================================================================
%%  SECTION 1 — INTRODUCTION
%% ===================================================================
\section{Introduction}

\IEEEPARstart{R}{eliability} processes are required in power system
operations due to a
fundamental market failure in the provision of reliability. To address
this, system operators routinely commit generation resources outside of
the energy market through reliability-driven commitment actions.
While effective in preserving system feasibility, empirical evidence
and operational experience indicate that the heuristics used to trigger
these commitments are often overly conservative, leading to systematic
over-commitment and distortion of price signals in both the day-ahead
and real-time markets.

Motivated by this gap, this work proposes a robust optimization
framework to identify the minimal set of generation commitment decisions
that are strictly necessary prior to the next day-ahead market clearing
to ensure system reliability. Rather than optimizing all commitment and
dispatch decisions under uncertainty, the proposed methodology focuses
exclusively on commitment decisions that cannot be deferred without
risking infeasibility under adverse realizations of load and renewable
generation.

System reliability is enforced with respect to a worst-case forecast
model of net demand, capturing uncertainty in both load and renewable
production. Commitment decisions are selected so as to guarantee
feasibility for all realizations within this uncertainty set, while
minimizing the extent of out-of-market intervention. In this sense, the
resulting commitments may be interpreted as the least-distortionary
reliability actions, providing a principled defense of operator
interventions while preserving market signals to the greatest extent
possible.

\subsection{Contributions}

This paper makes the following contributions.

First, we develop a notification-gated formulation for reliability unit
commitment that explicitly models the interaction between physical
generator lead times and the sequencing of market and operational
decision processes. The proposed framework identifies the subset of
commitment decisions that cannot be deferred to subsequent market stages
and restricts out-of-market interventions to this non-deferrable
decision set.

Second, we introduce a minimal-intervention reliability objective that
minimizes only startup actions and associated no-load costs induced by
non-deferrable commitments, while enforcing worst-case dispatch
feasibility under ellipsoidal uncertainty through linear decision rules.
The resulting mixed-integer second-order cone program jointly optimizes
commitment and robust dispatch, decoupling reliability enforcement from
full economic re-optimization while preserving market-based outcomes.

Third, we demonstrate that the proposed formulation generalizes
naturally across operational horizons, including multi-day reliability
assessment and day-ahead reliability unit commitment, enabling
coordinated treatment of long-lead thermal resources and near-term
reliability adjustments within a unified optimization framework.

Fourth, we conduct a comprehensive numerical evaluation comparing the
proposed approach against standard reliability unit commitment and
integrated market-reliability formulations. Results show significant
reductions in early commitment interventions and redispatch magnitude
without compromising system reliability.


%% ===================================================================
%%  SECTION 2 — LITERATURE REVIEW
%% ===================================================================
\section{Literature Review}

\subsection{Reliability Commitment in Practice}

Every North American RTO operates a reliability unit commitment
process following the day-ahead market clearing.
The Southwest Power Pool (SPP) runs a DARUC with a 48-hour
look-ahead that may commit additional resources beyond the
day-ahead market solution~[cite SPP market protocols].
The Midcontinent ISO (MISO) executes a similar reliability
assessment commitment (RAC) process to ensure resource adequacy
over the next operating day~[cite MISO BPM-002].
CAISO's residual unit commitment (RUC) procures capacity to
close any gap between day-ahead scheduled supply and the
operator's demand forecast~[cite CAISO BPM].
PJM, NYISO, and ERCOT employ analogous processes under varying
nomenclature~[cite PJM manual 11; NYISO OATT; ERCOT protocols].

Despite the ubiquity of these processes, their mathematical
formulations are largely proprietary and vary considerably across
operators.
FERC Order~2222 and subsequent technical conferences have drawn
attention to the market impact of reliability commitments, noting
that out-of-market actions can suppress energy prices and
create uplift charges~[cite FERC].
The academic literature has given comparatively little attention
to the formulation of the reliability commitment problem itself,
typically treating it as a standard unit commitment solved
after the market clears.
The sequential nature of the decision process --- and in
particular the question of \emph{which} decisions must be made
at each stage --- has received almost no formal treatment.

\subsection{Robust and Stochastic Unit Commitment}

The robust optimization approach to unit commitment was
pioneered by Bertsimas et~al.~[cite Bertsimas 2012], who
introduced a polyhedral uncertainty set for net load and solved
the two-stage problem via Benders decomposition.
Subsequent work by Jiang et~al.~[cite Jiang 2012] developed
column-and-constraint generation methods that avoid the need for
dual-based decomposition and demonstrated improved computational
performance.
Zhao and Guan~[cite Zhao Guan 2013] extended the framework to
multi-stage settings and established finite convergence results
for the CCG procedure.

Stochastic unit commitment provides an alternative treatment of
uncertainty through scenario-based formulations.
Papavasiliou et~al.~[cite Papavasiliou 2011] demonstrated the
value of stochastic optimization for wind integration, and
Zheng et~al.~[cite Zheng 2015] developed decomposition methods
for large-scale stochastic UC problems.
Distributionally robust approaches, which optimize over
ambiguity sets of probability distributions, have been explored
by Xiong et~al.~[cite Xiong 2017] and others.

Linear decision rules (LDR) provide a tractable approximation of
the recourse policy in robust and adaptive optimization.
Ben-Tal et~al.~[cite Ben-Tal 2004] established the affine
adjustable robust counterpart framework, and Lorca and
Sun~[cite Lorca Sun 2016] applied LDR to unit commitment with
ellipsoidal uncertainty sets, obtaining second-order cone
reformulations.
The present work builds on the LDR-based robust UC framework
developed in~[paper~2], which introduces learned covariance
estimation and conformal prediction for uncertainty set
calibration.

A common feature of the existing robust UC literature is that
the formulation scope encompasses \emph{all} commitment decisions
over the planning horizon.
The contribution of this paper is orthogonal to the choice of
uncertainty model: rather than proposing a new uncertainty
representation, we introduce a \emph{decision scope} restriction
that limits the reliability optimization to non-deferrable
commitments, regardless of whether the underlying model is
deterministic, robust, or stochastic.

\subsection{Market Efficiency and Out-of-Market Actions}

Out-of-market reliability actions have well-documented effects on
market efficiency.
Reliability commitments that deviate from the market solution
create make-whole payments (uplift) to committed generators,
representing costs that are not reflected in locational marginal
prices~[cite Hogan 2003; O'Neill 2005].
The convex hull pricing literature~[cite Gribik 2007;
Schiro 2016] has established that pricing in the presence of
non-convexities (including commitment decisions) cannot
simultaneously achieve revenue adequacy and cost recovery without
some form of side payment.

Reducing the number of out-of-market commitments directly
reduces the magnitude of uplift and improves the extent to which
energy prices reflect the true marginal cost of supply.
Andrianesis et~al.~[cite Andrianesis 2020] quantify the
relationship between commitment interventions and uplift in
ISO-NE, finding that reliability commitments account for a
substantial fraction of total uplift payments.
Ela et~al.~[cite Ela 2011] analyze the operational and market
impacts of wind integration on commitment processes, noting that
conservative commitment practices designed to hedge against
renewable forecast errors contribute to over-commitment and
price distortion.

\subsection{Research Gap}

The existing literature treats the reliability commitment problem
as a monolithic optimization over the full planning horizon:
once the market clears, the reliability process re-optimizes all
remaining commitment decisions under its chosen uncertainty model.
No existing formulation distinguishes between commitment decisions
that must be finalized at the current stage and those that can be
deferred to subsequent processes.

This paper addresses this gap by introducing the gating set, which
uses physical generator notification times to partition the
decision space into deferrable and non-deferrable components.
The result is a principled restriction of the reliability
optimization scope that preserves market outcomes to the greatest
extent permitted by physical constraints.


%% ===================================================================
%%  SECTION 3 — THE RELIABILITY COMMITMENT PROBLEM
%% ===================================================================
\section{The Reliability Commitment Problem}

\subsection{Operational Timeline}

North American RTOs share a common sequential structure in their
commitment and dispatch processes, though specific timing and
nomenclature vary across operators.
For notational clarity, we refer to the \emph{Operating Day}~(OD)
as the day for which the operation plan is being generated.
The following sequence is representative of the general pattern;
specific timing parameters are given in
Section~\ref{sec:casestudy}.

The earliest commitment decisions are made through a
\emph{multi-day reliability assessment}~(MDRA), an advisory process
run several days before OD with a look-ahead horizon of four to
seven days.
The MDRA identifies long-lead resources --- typically large thermal
units requiring multi-day advance notice --- that must be committed
to ensure system adequacy.

The principal scheduling process is the \emph{Day-Ahead
Market}~(DAMKT), which clears on OD$-$1 and produces a binding
commitment and dispatch schedule for~OD.
Following the DAMKT, the \emph{Day-Ahead Reliability Unit
Commitment}~(DARUC) reviews the market solution and commits
additional resources as needed to ensure reliability over a horizon
covering OD and, in many implementations, OD$+$1.

During the operating day itself, \emph{Intra-Day RUC}~(ID-RUC)
processes run on a rolling basis, making near-term commitment
adjustments as updated forecasts become available.

The complete sequence --- MDRA, DAMKT, DARUC, ID-RUC --- forms a
nested hierarchy in which each successive process operates over a
shorter horizon with more current information, and may revise or
supplement the decisions of its predecessor.

\subsection{Unit Commitment Formulation}

% State the canonical UC formulation ONCE. Both DAMKT and standard
% DARUC share these constraints; they differ only in objective.

The following formulation captures the core unit commitment
constraints common to both the day-ahead market and reliability
processes. Let $\mathcal{I}$ index generating resources,
$\mathcal{T} = \{1,\dots,T\}$ index operating periods,
$\mathcal{N}$ index buses, and $\mathcal{L}$ index transmission
lines.

\begin{subequations}\label{eq:uc_base}
\begin{align}
%--- power balance ---
&\sum_{i} p_{i,t} + s_t^{p} = \sum_n d_{n,t},
  \quad \forall t, \label{eq:uc_bal} \\
%--- transmission ---
&-\bar{F}^l \leq \sum_n S_n^l
  \Big(\sum_{i\in\mathcal{I}_n} p_{i,t} - d_{n,t}\Big)
  \leq \bar{F}^l,
  \quad \forall l,t, \label{eq:uc_flow} \\
%--- block dispatch ---
&p_{i,t,b} \leq \bar{P}_{i,b}\, u_{i,t},
  \;\; \forall i,t,b, \quad
  p_{i,t} = \textstyle\sum_b p_{i,t,b},
  \;\; \forall i,t, \label{eq:uc_block} \\
%--- renewable limit ---
&p_{j,t} \leq \bar{p}^{\text{fcst}}_{j,t},
  \quad \forall j \in \mathcal{I}^{\text{ren}},\; t, \label{eq:uc_ren} \\
%--- capacity ---
&\underline{P}_i\, u_{i,t} \leq p_{i,t}
  \leq \bar{P}_i\, u_{i,t},
  \quad \forall i,t, \label{eq:uc_cap} \\
%--- ramp up ---
&p_{i,t} - p_{i,t-1} \leq R_i^U (u_{i,t-1} + v_{i,t}),
  \quad \forall i,t, \label{eq:uc_ru} \\
%--- ramp down ---
&p_{i,t-1} - p_{i,t} \leq R_i^D (u_{i,t} + w_{i,t}),
  \quad \forall i,t, \label{eq:uc_rd} \\
%--- logic ---
&u_{i,t} - u_{i,t-1} = v_{i,t} - w_{i,t},
  \quad \forall i,t, \label{eq:uc_logic} \\
%--- MUT ---
&\sum_{t'=t-\text{MUT}_i+1}^{t} v_{i,t'} \leq u_{i,t},
  \quad \forall i,\; t \geq \text{MUT}_i, \label{eq:uc_mut} \\
%--- MDT ---
&\sum_{t'=t-\text{MDT}_i+1}^{t} w_{i,t'} \leq 1 - u_{i,t},
  \quad \forall i,\; t \geq \text{MDT}_i, \label{eq:uc_mdt} \\
%--- initial conditions ---
&\sum_{t'=1}^{UT_i} u_{i,t'} = UT_i, \quad
  \sum_{t'=1}^{DT_i} u_{i,t'} = 0,
  \quad \forall i, \label{eq:uc_init} \\
%--- domains ---
&\mathbf{p},\, \mathbf{s} \geq 0, \quad
  \mathbf{u},\, \mathbf{v},\, \mathbf{w} \in \{0,1\}.
  \label{eq:uc_domain}
\end{align}
\end{subequations}

We denote this shared constraint set as $\mathcal{X}^{\text{UC}}$.

\subsection{Day-Ahead Market (DAMKT)}

The DAMKT minimizes total production cost including energy,
no-load, startup, and shutdown costs over all generators and periods:
\begin{align}\label{eq:damkt_obj}
\min_{\mathbf{p},\mathbf{u},\mathbf{v},\mathbf{w},\mathbf{s}}
\quad &\sum_{i,t} \Big[
  C_i^{NL} u_{i,t} + C_i^{SU} v_{i,t} + C_i^{SD} w_{i,t}
  + \sum_b C_{i,b}^p\, p_{i,t,b}
\Big] \nonumber\\
&+ \sum_t M^p s_t^p \\
\text{s.t.} \quad
& \eqref{eq:uc_base}. \nonumber
\end{align}

\subsection{Standard Day-Ahead RUC (DARUC)}

The standard DARUC uses the same constraint set
$\mathcal{X}^{\text{UC}}$ but removes energy costs from the
objective, retaining only commitment-related costs:
\begin{align}\label{eq:daruc_obj}
\min_{\mathbf{p},\mathbf{u},\mathbf{v},\mathbf{w},\mathbf{s}}
\quad &\sum_{i,t} \Big[
  C_i^{NL} u_{i,t} + C_i^{SU} v_{i,t} + C_i^{SD} w_{i,t}
\Big]
+ \sum_t M^p s_t^p \\
\text{s.t.} \quad
& \eqref{eq:uc_base}. \nonumber
\end{align}

In practice, the DARUC typically inherits the DAMKT commitment
schedule as a lower bound, permitting additional commitments but
prohibiting de-commitment of market-cleared resources.
The specific implementation varies across operators; some fix the
market schedule entirely, while others allow limited modifications.
This no-decommitment convention is formalized in
Section~\ref{sec:decomp}.

\subsection{Why Standard DARUC Over-Commits}\label{sec:overcommit}

The standard DARUC formulation optimizes commitment decisions
$u_{i,t}$ for every generator~$i$ across every period~$t$ in the
planning horizon.
However, the vast majority of these decisions will be revisited ---
and potentially revised --- by subsequent processes: the next DAMKT
clearing for periods beyond~OD, and intra-day RUC for near-term
adjustments.
The DARUC thus solves an optimization problem whose scope extends
well beyond the decisions that are genuinely its to make.

This scope mismatch creates a systematic bias toward early
commitment.
Because the DARUC minimizes total commitment cost over the full
horizon, it may find it optimal to start a unit in an early period
to avoid a more expensive startup in a later period --- even when
the later decision falls within the purview of a subsequent process
that will have access to updated forecasts.
The DARUC has no mechanism to recognize that such deferred decisions
carry no immediate cost: from its perspective, the entire horizon
is a single optimization to be solved jointly.

The result is commitment decisions that are feasible but premature.
Units are started hours or days before they are operationally
required, displacing resources that were optimally scheduled by the
market and distorting the price signals that the market was designed
to produce.

The fundamental issue is not that the DARUC commits the wrong units
in an absolute sense, but that it commits units whose startup
decisions could have been deferred without jeopardizing reliability.
The formulation lacks a notion of \emph{deferrability}: it treats
every commitment decision as equally urgent, regardless of whether
the physical lead time of the generator demands an immediate
decision or permits deferral to a later stage.
Section~\ref{sec:gating} introduces the gating set to formalize
this distinction.


%% ===================================================================
%%  SECTION 4 — NOTIFICATION-GATED RUC (THE CONTRIBUTION)
%% ===================================================================
\section{Notification-Gated Reliability Commitment}

The proposed formulation jointly optimizes commitment decisions and
robust dispatch in a single mixed-integer second-order cone program
(MISOCP).
The gating set restricts the reliability objective to non-deferrable
interventions, while linear decision rules and ellipsoidal uncertainty
sets enforce worst-case dispatch feasibility across all constraint
families.

\subsection{Gating Sets}\label{sec:gating}

The notification (lead) time of a generating unit determines which
commitment decisions must be finalized at the current decision stage
and which can be deferred to subsequent market or operational
processes.
Commitment decisions should be made binding only when deferral would
violate physical feasibility.

Let decision process $k$ represent the current reliability or market
commitment stage (e.g., MDRA, DA-RUC, ID-RUC), and let
$t^{\text{next}}$ denote the first operating period whose commitment
decisions can be modified by the next available process.
Let $t \in \mathcal{T}^k = \{1,\dots,T\}$ index the physical
operating periods in the look-ahead horizon considered at stage~$k$.

For generating resource~$i$ with startup notification time
$L_i^{SU}$, a startup decision affecting operating period~$t$ must
be issued at the current stage if it cannot be deferred to the next
process.
This condition defines the \emph{gating set}:
\begin{equation}\label{eq:gating_set}
    \mathcal{T}_{i}^{\text{gate}} \;:=\;
    \bigl\{\, t \in \mathcal{T}^k \;\big|\;
            t - L_i^{SU} < t^{\text{next}} \bigr\}.
\end{equation}

Periods in $\mathcal{T}_{i}^{\text{gate}}$ correspond to startup
decisions that must be initiated before the next commitment process
in order to preserve feasibility.
Startup decisions outside this set may be deferred and are excluded
from the current reliability optimization.

Because $t - L_i^{SU}$ is increasing in~$t$, the gating set is a
contiguous block of early periods: if period~$t$ is gated, then every
earlier period $t' < t$ is also gated.
Generators with longer notification times have larger gating sets,
reflecting the greater advance notice required to start them.

The gating set restricts \emph{startup} decisions ($v_{i,t}$), not
commitment status ($u_{i,t}$).
A unit already committed by the market does not require a new startup
decision and is therefore not subject to gating.

\subsection{Market-Reliability Decomposition}\label{sec:decomp}

Let $(\mathbf{u}^m, \mathbf{v}^m, \mathbf{w}^m)$ denote the
commitment schedule from the preceding DAMKT solution.
The reliability process operates on the \emph{total} commitment
variables $(u_{i,t}, v_{i,t}, w_{i,t})$, which are constrained to
preserve all market-decided commitments:
\begin{equation}\label{eq:nodecommit}
    u_{i,t} \;\geq\; u_{i,t}^m, \qquad \forall\, i,t.
\end{equation}
The startup and shutdown variables $v_{i,t}$ and $w_{i,t}$ are
determined by the logic constraint~\eqref{eq:uc_logic} applied to the
total commitment~$u_{i,t}$ and are not independently decomposed.

The \emph{incremental commitment} is defined as
\begin{equation}\label{eq:uprime}
    u'_{i,t} \;:=\; u_{i,t} - u_{i,t}^m \;\in\; \{0,1\},
    \qquad \forall\, i,t,
\end{equation}
where $u'_{i,t} = 1$ indicates a reliability intervention: a unit
brought online beyond the market solution.

A reliability startup at period~$t$ may \emph{absorb} a market
startup at a later period: if the market planned to start unit~$i$ at
period~$t{+}2$, but the reliability process starts it at~$t$, the
unit is already online at~$t{+}2$ and the market startup is
superseded.
This interaction is handled automatically by the logic
constraint~\eqref{eq:uc_logic} operating on total~$u$.

\subsection{Uncertainty Model}\label{sec:uncertainty}

Let $\mathbf{r} \in \mathbb{R}^K$ represent deviations of renewable
generation from the mean forecast, where $K$ is the number of
uncertain renewable sources.
The uncertainty set is an ellipsoid parameterized by a positive
definite covariance matrix
$\boldsymbol{\Sigma} \in \mathbb{R}^{K \times K}$ and a scalar
budget~$\rho > 0$:
\begin{equation}\label{eq:uset}
    \mathcal{U} \;:=\;
    \bigl\{\, \mathbf{r} \in \mathbb{R}^K \;\big|\;
    \mathbf{r}^\top \boldsymbol{\Sigma}^{-1} \mathbf{r}
    \leq \rho^2 \bigr\}.
\end{equation}
The construction of $\boldsymbol{\Sigma}$ and calibration of~$\rho$
via conformal prediction are detailed in~[paper~2]; here we treat
$(\boldsymbol{\Sigma}, \rho)$ as given inputs.

Dispatch is modeled as an affine function of the uncertainty
realization via linear decision rules (LDR):
\begin{equation}\label{eq:ldr}
    p_{i,t}(\mathbf{r})
    \;=\; p_{i,t}^0
    \;+\; \mathbf{z}_{i,t}^\top \mathbf{r},
    \qquad \forall\, i,t,
\end{equation}
where $p_{i,t}^0$ is the nominal dispatch and
$\mathbf{z}_{i,t} \in \mathbb{R}^K$ are the LDR coefficients
governing the response of generator~$i$ at period~$t$ to the
uncertainty realization.

For a generic constraint of the form
$\mathbf{z}_{i,t}^\top \mathbf{r} \leq b_{i,t}$ that must hold for
every $\mathbf{r} \in \mathcal{U}$, the worst case yields the
second-order cone (SOC) condition
\begin{equation}\label{eq:soc_generic}
    \rho\,
    \bigl\lVert
      \boldsymbol{\Sigma}^{1/2} \mathbf{z}_{i,t}
    \bigr\rVert_2
    \;\leq\; b_{i,t}.
\end{equation}


\subsection{Gated Robust LD-RUC Formulation}\label{sec:ldruc}

The notification-gated LD-RUC is formulated as a single MISOCP that
simultaneously determines the minimum-cost set of non-deferrable
reliability commitments and a dispatch policy that is feasible under
all uncertainty realizations.
The commitment variables $(\mathbf{u}, \mathbf{v}, \mathbf{w})$ are
binary; the nominal dispatch~$\mathbf{p}^0$ and LDR
coefficients~$\mathbf{Z}$ are continuous.

\begin{subequations}\label{eq:ldruc}
\begin{align}
\min_{\substack{
  \mathbf{u},\mathbf{v},\mathbf{w},\\
  \mathbf{p}^0,\mathbf{Z},\mathbf{s}}} \quad
& \sum_{i \in \mathcal{I}}
  \sum_{t \in \mathcal{T}_i^{\text{gate}}}
  \!\Big[\,
    C_i^{SU}\, v_{i,t}
    \;+\; C_i^{NL}\, u'_{i,t}
  \Big]
  \;+\; M \!\sum_{t} s_t
  \label{eq:ldruc_obj} \\[4pt]
\text{s.t.}\quad
%--- no decommitment ---
& u_{i,t} \geq u_{i,t}^m,
  \quad \forall\, i,t,
  \label{eq:ldruc_lock} \\
%--- incremental commitment ---
& u'_{i,t} = u_{i,t} - u_{i,t}^m,
  \quad \forall\, i,t,
  \label{eq:ldruc_uprime} \\
%--- nominal power balance ---
& \sum_{i} p_{i,t}^0 + s_t = \sum_n d_{n,t},
  \quad \forall\, t,
  \label{eq:ldruc_bal_nom} \\
%--- recourse balance ---
& \sum_{i} z_{i,t,k} = 0,
  \quad \forall\, k,t,
  \label{eq:ldruc_bal_rec} \\
%--- capacity upper (SOC) ---
& p_{i,t}^0
  + \rho \lVert \boldsymbol{\Sigma}^{1/2} \mathbf{z}_{i,t} \rVert_2
  \leq \bar{P}_i\, u_{i,t},
  \quad \forall\, i,t,
  \label{eq:ldruc_cap_up} \\
%--- capacity lower (SOC) ---
& \underline{P}_i\, u_{i,t}
  \leq p_{i,t}^0
  - \rho \lVert \boldsymbol{\Sigma}^{1/2} \mathbf{z}_{i,t} \rVert_2,
  \quad \forall\, i,t,
  \label{eq:ldruc_cap_lo} \\
%--- ramp up (SOC) ---
& (p_{i,t}^0 - p_{i,t-1}^0)
  + \rho \lVert \boldsymbol{\Sigma}^{1/2}
    (\mathbf{z}_{i,t} - \mathbf{z}_{i,t-1}) \rVert_2
  \nonumber\\
& \qquad\qquad
  \leq R_i^U (u_{i,t-1} + v_{i,t}),
  \quad \forall\, i,t,
  \label{eq:ldruc_ramp_up} \\
%--- ramp down (SOC) ---
& (p_{i,t-1}^0 - p_{i,t}^0)
  + \rho \lVert \boldsymbol{\Sigma}^{1/2}
    (\mathbf{z}_{i,t-1} - \mathbf{z}_{i,t}) \rVert_2
  \nonumber\\
& \qquad\qquad
  \leq R_i^D (u_{i,t} + w_{i,t}),
  \quad \forall\, i,t,
  \label{eq:ldruc_ramp_dn} \\
%--- transmission (SOC) ---
& \Big|\sum_n S_n^l
  \Big(\sum_{i \in \mathcal{I}_n} p_{i,t}^0 - d_{n,t}\Big)\Big|
  + \rho \Big\lVert \boldsymbol{\Sigma}^{1/2}
    \sum_{i \in \mathcal{I}_n} S_n^l\, \mathbf{z}_{i,t}
  \Big\rVert_2
  \nonumber\\
& \qquad\qquad
  \leq \bar{F}^l,
  \quad \forall\, l,t,
  \label{eq:ldruc_flow} \\
%--- renewable limit (SOC) ---
& p_{j,t}^0
  + \rho \lVert \boldsymbol{\Sigma}^{1/2}
    (\mathbf{z}_{j,t} - \mathbf{e}_{j,t}) \rVert_2
  \leq \mu_{j,t},
  \nonumber\\
& \qquad\qquad
  \forall\, j \in \mathcal{I}^{\text{ren}},\; t,
  \label{eq:ldruc_ren} \\
%--- logic ---
& u_{i,t} - u_{i,t-1} = v_{i,t} - w_{i,t},
  \quad \forall\, i,t,
  \label{eq:ldruc_logic} \\
%--- MUT/MDT ---
& \text{MUT/MDT constraints}~\eqref{eq:uc_mut}\text{--}\eqref{eq:uc_mdt},
  \label{eq:ldruc_mutmdt} \\
%--- initial conditions ---
& \text{initial conditions}~\eqref{eq:uc_init},
  \label{eq:ldruc_init} \\
%--- domains ---
& \mathbf{p}^0 \geq 0,\; \mathbf{s} \geq 0, \quad
  \mathbf{u},\, \mathbf{v},\, \mathbf{w} \in \{0,1\}.
  \label{eq:ldruc_domain}
\end{align}
\end{subequations}

Here $\mathbf{e}_{j,t} \in \mathbb{R}^K$ is the unit vector
selecting the uncertainty dimension corresponding to renewable
source~$j$ at period~$t$, so that the available renewable output
under realization~$\mathbf{r}$ is $\mu_{j,t} + \mathbf{e}_{j,t}^\top
\mathbf{r}$.
Constraint~\eqref{eq:ldruc_bal_rec} ensures that the aggregate
dispatch response to any realization sums to zero: when renewable
output deviates, the remaining fleet absorbs the imbalance.

\emph{Structure.}
Formulation~\eqref{eq:ldruc} is a mixed-integer second-order cone
program (MISOCP).
The integer variables arise from the commitment decisions
$(\mathbf{u}, \mathbf{v}, \mathbf{w})$; the continuous variables
$(\mathbf{p}^0, \mathbf{Z})$ determine the robust dispatch policy.
The SOC constraints~\eqref{eq:ldruc_cap_up}--\eqref{eq:ldruc_ren}
couple the commitment and dispatch decisions: when a unit is
de-committed ($u_{i,t} = 0$), the corresponding capacity and ramp
bounds force $p_{i,t}^0 = 0$ and $\mathbf{z}_{i,t} = \mathbf{0}$.
Modern MISOCP solvers (e.g., Gurobi) handle this structure via
branch-and-bound with SOCP relaxations at each node.

\emph{Objective-only gating.}
The objective~\eqref{eq:ldruc_obj} penalizes startup and
incremental no-load costs only within each generator's gating window
$\mathcal{T}_i^{\text{gate}}$.
However, the \emph{constraints} apply over the full horizon: the
optimizer is free to add non-gated commitments when forced by
minimum-up-time propagation or system feasibility, but receives no
objective incentive to do so.
This design ensures that the formulation accounts for the physical
consequences of gated decisions (e.g., a gated startup at period~$t$
commits the unit through $t + \text{MUT}_i - 1$) without penalizing
those consequences beyond the gating window.

\emph{Windowed no-load variant.}
A gated startup at period~$t$ commits unit~$i$ for at least
$\text{MUT}_i$ periods, some of which may fall beyond the gating
window.
A refined objective attributes no-load cost over the irrevocable
commitment window rather than period-by-period:
\begin{equation}\label{eq:obj_window}
\sum_{i \in \mathcal{I}}
\sum_{t \in \mathcal{T}_i^{\text{gate}}}
\!\left[
  C_i^{SU}\, v_{i,t}
  \;+\;
  \sum_{\tau=t}^{\min\{t+\text{MUT}_i-1,\,T\}}
    \!\! C_i^{NL}\, v'_{i,t}
\right]
+ M \!\sum_t s_t,
\end{equation}
where $v'_{i,t} := v_{i,t}\,(1 - v_{i,t}^m)$ isolates reliability
startups from market startups.

\emph{Relationship to existing formulations.}
The LD-RUC formulation~\eqref{eq:ldruc} differs from the standard
robust DARUC in two respects:
\begin{enumerate}
\item \emph{Gated objective.} Only commitment costs within the
  gating window contribute to the objective. The standard DARUC
  penalizes commitment costs over the full horizon.
\item \emph{No-decommitment floor.} The market commitment
  $\mathbf{u}^m$ is preserved as a lower bound via~\eqref{eq:ldruc_lock}.
  While this constraint also appears in standard DARUC implementations,
  its interaction with the gated objective is what produces the
  least-interventionist behavior.
\end{enumerate}
The robust dispatch constraints
\eqref{eq:ldruc_cap_up}--\eqref{eq:ldruc_ren} are identical to
those used in the adaptive robust UC (ARUC) framework
of~[paper~2].
The contribution here is the decision scope restriction, not the
uncertainty model.


\subsection{Summary}\label{sec:formulation_summary}

The complete LD-RUC procedure is:
\begin{enumerate}
  \item Solve the DAMKT~\eqref{eq:damkt_obj} to obtain the market
    commitment $\mathbf{u}^m$.
  \item Compute gating sets $\mathcal{T}_i^{\text{gate}}$ from
    notification times and $t^{\text{next}}$
    via~\eqref{eq:gating_set}.
  \item Solve the gated robust LD-RUC~\eqref{eq:ldruc}: a single
    MISOCP that jointly determines incremental commitments and a
    worst-case feasible dispatch policy.
\end{enumerate}

The formulation requires no iterative procedure: the MISOCP is
solved once and the resulting commitment schedule is guaranteed to be
robust-feasible by construction.
This is a direct consequence of jointly optimizing commitment and
dispatch under the SOC robust counterpart --- the same approach used
by the ARUC framework, but with the objective restricted to the
gating window.


%% ===================================================================
%%  SECTION 5 — GENERALIZATION ACROSS DECISION STAGES
%% ===================================================================
\section{Generalization Across Decision Stages}\label{sec:general}

The gating set~\eqref{eq:gating_set} is parameterized by the
decision stage~$k$ (through $t^{\text{next}}$) and the generator
notification times~$L_i^{SU}$.
By varying these parameters, the LD-RUC framework instantiates
naturally across the operational planning sequence.

\subsection{Multi-Day Reliability Assessment (MDRA)}

The MDRA is an advisory process run several days before the operating
day, with a horizon of four to seven days.
At this stage, $t^{\text{next}}$ corresponds to the first period of
OD$-$1 (the day before the operating day), when the DAMKT will
clear.
Only generators with very long notification times --- those whose
startup lead time exceeds the interval between the MDRA and the
DAMKT --- have non-empty gating sets.
In practice, this restricts the MDRA's binding decisions to large
coal, nuclear, and combined-cycle units requiring multi-day advance
notice.

\subsection{Day-Ahead RUC (DA-RUC)}

The DA-RUC runs after the DAMKT clears for the operating day.
Here $t^{\text{next}}$ is the first period of OD$+$1 (the next
operating day), since the next DAMKT will cover that horizon.
Generators with notification times shorter than 24~hours have gating
sets that cover only the operating day itself; those with longer lead
times may also gate periods on OD$+$1.
This is the primary application of the LD-RUC framework developed in
Section~\ref{sec:ldruc}.

\subsection{Intra-Day RUC (ID-RUC)}

Intra-day reliability processes run on a rolling basis during the
operating day.
At each invocation, $t^{\text{next}}$ is the earliest period that can
still be modified by the next intra-day run (typically the current
period plus the process cycle time).
The gating sets shrink as the operating day progresses and more
decisions become irrevocable, leaving only fast-start resources with
non-empty gating windows.

\subsection{Relationship to Standard DARUC}

The standard DARUC formulation~\eqref{eq:daruc_obj} can be viewed as
the limiting case of the gating framework in which
$\mathcal{T}_i^{\text{gate}} = \mathcal{T}$ for all~$i$ --- that is,
every commitment decision is treated as non-deferrable.
This occurs when $t^{\text{next}} \to \infty$ (no subsequent process
exists) or equivalently when $L_i^{SU} \to \infty$ for all
generators.

Conversely, the pure market outcome (no reliability intervention)
corresponds to $\mathcal{T}_i^{\text{gate}} = \varnothing$ for
all~$i$, which occurs when $L_i^{SU} = 0$ (all units can be started
instantaneously) or $t^{\text{next}} = 1$ (the next process can
modify every period).
The gating framework thus spans a continuum between no intervention
and full re-optimization, with the physical notification times
determining the appropriate point on this spectrum.


%% ===================================================================
%%  SECTION 6 — CASE STUDY
%% ===================================================================
\section{Case Study}\label{sec:casestudy}

\subsection{Setup}

The proposed LD-RUC formulation is evaluated on the RTS-GMLC test
case, a three-area system with 73~generating units (thermal, wind,
solar, and hydro) and 120~transmission lines.
Wind forecast data are drawn from the Southwest Power Pool (SPP)
constellation forecast dataset, scaled to the RTS-GMLC wind fleet
capacities.
The ellipsoidal uncertainty set parameters
$(\boldsymbol{\Sigma}, \rho)$ are constructed using the learned
covariance and conformal prediction methodology developed
in~[paper~2].

Notification times $L_i^{SU}$ are assigned to the generator fleet
based on fuel type and capacity:
% TODO: Table of notification time assignments.
% Coal/nuclear: 24-48 hours
% Combined cycle: 8-12 hours
% Combustion turbine: 1-4 hours
% Wind/solar/hydro: 0 hours (no startup lead time)

The DA-RUC decision stage is used as the primary test case, with
$t^{\text{next}}$ set to the first period of OD$+$1 (period~25 for a
48-hour horizon starting at OD hour~1).
All experiments are solved using Gurobi~11 on [hardware description].

\subsection{Standard DARUC vs.\ LD-RUC}\label{sec:results_daruc}

% TODO: Primary comparison results.
% Expected structure:
% - Table: number of reliability commitments (startups) by fuel type
%   for standard DARUC vs LD-RUC
% - Figure: commitment heatmap showing which generators are committed
%   differently between the two approaches
% - Key finding: long-lead units (coal, CC) have identical commitments
%   since they are gated in both formulations; short-lead units (CT)
%   differ significantly
% - Total commitment cost comparison
% - Feasibility verification: confirm no load shedding (s = 0)
%   in both cases

\subsection{Comparison with Integrated Robust UC}
\label{sec:results_aruc}

% TODO: Compare LD-RUC against the integrated ARUC from paper 2.
% - ARUC jointly optimizes all commitment and dispatch decisions
%   under uncertainty (no gating, no decomposition)
% - Metrics: total cost, commitment count, wind curtailment,
%   computation time
% - Key point: LD-RUC produces fewer out-of-market interventions
%   than ARUC by construction, at the potential cost of higher
%   total system cost
% - The cost gap between LD-RUC and ARUC represents the price of
%   minimal intervention: how much more expensive is the
%   least-distortionary solution?

\subsection{Sensitivity to Notification Times}
\label{sec:results_sensitivity}

% TODO: This section demonstrates the gating set's value directly.
% - Sweep L_i^SU for the fleet (or sweep t^next)
% - For each setting: report number of gated decisions, number of
%   reliability commitments, total cost
% - Figure: commitments and cost vs notification time or t^next
% - Limiting cases as validation:
%   L_i -> 0: no gated decisions, zero reliability cost
%   L_i -> infinity: all decisions gated, recovers standard DARUC
% - Intermediate values: demonstrate smooth transition and
%   quantify the marginal value of each hour of notification time


%% ===================================================================
%%  SECTION 7 — CONCLUSION
%% ===================================================================
\section{Conclusion}\label{sec:conclusion}

This paper introduced a notification-gated formulation for
reliability unit commitment that restricts out-of-market interventions
to the minimal set of commitment decisions that cannot be deferred to
subsequent market or operational processes.
The gating set, derived from physical generator notification times and
the timing of sequential decision stages, provides a principled
criterion for distinguishing non-deferrable reliability actions from
decisions that can be left to the market.

The formulation is a mixed-integer second-order cone program that
jointly optimizes incremental commitment decisions and robust dispatch
via linear decision rules.
By restricting the reliability objective to the gating window while
enforcing worst-case feasibility over the full horizon, the LD-RUC
produces the least-interventionist commitment schedule that guarantees
system reliability under the specified uncertainty set.

The gating framework generalizes across operational horizons.
Multi-day reliability assessment, day-ahead RUC, and intra-day RUC
are all instances of the same formulation with different
parameterizations of the gating set.
Standard DARUC is recovered as a limiting case in which all decisions
are treated as non-deferrable.

% TODO: Summarize numerical findings once case study is complete.
% Expected: significant reduction in reliability commitments,
% LD-RUC commitment schedule is robust-feasible by construction,
% computation times comparable to standard DARUC.

Future work includes analysis of market efficiency impacts (LMP
changes and uplift costs resulting from reduced reliability
interventions), extension to chance-constrained and distributionally
robust uncertainty models, and application to multi-settlement market
structures where the interaction between reliability processes and
real-time pricing is most consequential.


%% ===================================================================
%%  BACK MATTER
%% ===================================================================

\appendices
\section{Nomenclature}\label{app:nomenclature}

\begin{table}[h]
\renewcommand{\arraystretch}{1.15}
\centering
\begin{tabular}{ll}
\hline
\textbf{Symbol} & \textbf{Description} \\
\hline
\multicolumn{2}{l}{\emph{Sets and indices}} \\
$\mathcal{I}$, $i$ & Generating resources \\
$\mathcal{I}^{\text{ren}}$, $j$ & Renewable generators \\
$\mathcal{I}_n$ & Generators at bus $n$ \\
$\mathcal{N}$, $n$ & Buses \\
$\mathcal{L}$, $l$ & Transmission lines \\
$\mathcal{T}$, $t$ & Operating periods \\
$\mathcal{T}_i^{\text{gate}}$ & Gating set for generator $i$ \\
$\mathcal{U}$ & Uncertainty set \\
$K$, $k$ & Uncertainty dimensions \\
$b$ & Dispatch blocks \\
\hline
\multicolumn{2}{l}{\emph{Parameters}} \\
$C_i^{NL},C_i^{SU},C_i^{SD}$ & No-load, startup, shutdown cost \\
$C_{i,b}^p$ & Block energy cost \\
$\bar{P}_i$, $\underline{P}_i$ & Max/min capacity \\
$\bar{P}_{i,b}$ & Block capacity \\
$R_i^U$, $R_i^D$ & Ramp-up/down rate \\
$\text{MUT}_i$, $\text{MDT}_i$ & Min up/down time \\
$L_i^{SU}$ & Startup notification time \\
$d_{n,t}$ & Nodal demand \\
$\bar{F}^l$ & Line flow limit \\
$S_n^l$ & PTDF entry \\
$\mu_{j,t}$ & Mean renewable forecast \\
$\boldsymbol{\Sigma}$ & Uncertainty covariance matrix \\
$\rho$ & Uncertainty budget \\
$M$ & Slack penalty \\
$u_{i,t}^m$, $v_{i,t}^m$ & DAMKT commitment/startup \\
$t^{\text{next}}$ & First period modifiable by next process \\
\hline
\multicolumn{2}{l}{\emph{Variables}} \\
$u_{i,t}$ & Commitment status (binary) \\
$v_{i,t}$, $w_{i,t}$ & Startup/shutdown (binary) \\
$u'_{i,t}$ & Incremental commitment \\
$p_{i,t}^0$ & Nominal dispatch (MW) \\
$\mathbf{z}_{i,t}$ & LDR coefficients \\
$s_t$ & Power balance slack \\
\hline
\end{tabular}
\end{table}

\section*{Acknowledgment}
The authors would like to thank\ldots

%\bibliographystyle{IEEEtran}
%\bibliography{IEEEabrv,../bib/paper}

\end{document}
